\section{Step 45 Distinguish ordinary files from controlled read tests}\label{sha-s45}

Use a real file to compare the published digest with a host calculation. Separately use the CLI fixtures to force positive short reads, an unsuccessful open and a failure after a selected prefix. A real file with known bytes checks normal I/O; it does not by itself establish every error branch.

\subsection{Observe and compare}
Use recipes \texttt{s45-1}, \texttt{s45-2}, \texttt{s45-3}, \texttt{s45-4}, \texttt{s45-5}, \texttt{s45-6}. Prepare each in a separate source image from the current cumulative program. An internal checkpoint is inspected before its temporary or normal return; otherwise pause at \texttt{main\_done}.

Use the completed entries available at this point; keep their inputs and observation boundary explicit. The surrounding explanation identifies the state being compared.

\subsection{Expected observation}
Both read sizes give the same successful digest; errors close acquired descriptors and preserve output.

\Needspace{9\baselineskip}
\subsection{Check your understanding}
Does an ordinary small file force a read-error branch?

\begin{itemize}
\item \textbf{A} Yes, if its length is not a multiple of sixty-four.
\item \textbf{B} No; the error must be exercised deliberately.
\item \textbf{C} Only when the filename contains digits.
\end{itemize}
